%% 阶段报告：DeReFusion 复现 + 算子专业化诊断（中文版）
%% 编译：xelatex -interaction=nonstopmode stage_report_zh.tex （跑两遍）
%% 与 stage_report_en.tex 结构完全一致，两者独立可编译
\documentclass[preprint,12pt]{elsarticle}

\usepackage{amssymb}
\usepackage{amsmath}
\usepackage{booktabs}
\usepackage{graphicx}
\usepackage{float}
\usepackage{microtype}
\usepackage{xeCJK}
\setCJKmainfont{SimSun}[BoldFont=SimHei,ItalicFont=KaiTi]
\setCJKsansfont{SimHei}
\setCJKmonofont{FangSong}
\usepackage{hyperref}

\journal{阶段报告 —— 内部}

\begin{document}

\begin{frontmatter}

\title{DeReFusion 复现与算子专业化诊断：阶段性技术报告\tnoteref{t1}}
\tnotetext[t1]{复现对象：Hsieh \& Chen (2026), Applied Soft Computing 203, 116252；
并在其之上补充了"工况依赖与资产依赖的算子适用性"受控诊断。}

\author[nwpu]{Dongxu Jiang\corref{cor1}}
\ead{jiangdongxu04@gmail.com}
\cortext[cor1]{通讯作者。邮箱：jiangdongxu04@gmail.com.}
\affiliation[nwpu]{organization={},
            city={},
            country={}}

\begin{abstract}
本报告汇总截至 2026-09-12 已完成的工作：DeReFusion 的复现，以及围绕同一个问题的一串受控诊断——
输入的\textbf{结构状态}是否决定哪种算子更有效。首先在标普 500 及另外两个资产上复现论文的核心对比
（DLinear 线性基座、LSTM--Transformer 非线性残差分支、无参数加法融合，统一置于 RevIN 之下）。
随后按\textbf{因果}相对波动率对预测误差分层，发现非线性分支的边际价值是\textbf{资产依赖}而非普适：
标普 500 上显著为正且跨种子稳健，ETHUSD 方向一致但统计临界，BTCUSD 则显著\textbf{反号}。
容量敏感性实验（隐层宽度 $23\!\to\!128$，其余全部冻结）表明，早先"线性全面占优"的读数\textbf{部分来自
非线性算子的容量瓶颈}：容量充分时 ETHUSD 在 $36/36$ 个观测上偏好非线性，而 BTCUSD 在任何容量下都不偏好。
一个容量对齐的代理路由实验（12 个预先固定的结构状态、线性 vs.\ MLP、协议完全一致）未发现稳健的
\textbf{状态级}算子专业化。最后，对九个资产的探索性扫描显示交互效应\textbf{跨零}（六负三正，五个显著），
并识别出两个在 leave-one-out 下方向稳定的关联（已实现波动率与 $|\mathrm{ACF}_1|$）——仅作为
\textbf{候选}资产级结构规律报告。报告明确区分：证据\textbf{排除}了什么（样本级路由）、
\textbf{支持}了什么（资产级算子异质性）、以及\textbf{仍未解决}的是什么。
\end{abstract}

\begin{keyword}
复现 \sep 时间序列预测 \sep 算子专业化 \sep 结构状态 \sep 自适应融合 \sep 受控对比
\end{keyword}

\end{frontmatter}

\section{研究范围与问题}
\label{sec:intro}

本阶段严格遵循"证据优先"的流程：先复现参照结果，然后一次只问一个问题，且所有分组规则、阈值与算子定义
都在\textbf{看到结果之前}固定。报告围绕三个问题组织：

\begin{itemize}
\item \textbf{Q-A（复现）}：已发表的融合策略排序（加法融合 $>$ 线性基座 $>$ 可学习门控）能否在独立获取的数据上复现？
\item \textbf{Q-B（工况）}：非线性分支的边际贡献是否取决于输入窗口的波动状态？
\item \textbf{Q-C（专业化）}：输入的\textbf{结构状态}是否决定哪种算子更有效——即路由是否有实验依据？
\end{itemize}

所有实验在 CPU 上运行，采用冻结的 70/10/20 时序切分，报告逐样本配对 bootstrap 置信区间（4000 次重采样）
与配对胜率。观测到测试结果之后，未对任何超参数做过调整。

\section{参照对比的复现}
\label{sec:repro}

上游仓库的数据集不可再分发，因此重新抓取了十个资产的日线 OHLC（2016--2025）；行数与论文一致
（股指与个股 2513，外汇 2602，BTCUSD 3653，ETHUSD 2975）。表~\ref{tab:repro} 给出标普 500 在
$T\in\{1,24\}$ 两个步长上的对比（seed 2021）。

\begin{table}[!t]
\centering
\small
\caption{核心对比的复现（GSPC，seed 2021，归一化尺度）。}
\label{tab:repro}
\begin{tabular}{@{}lcccc@{}}
\toprule
模型 & $T{=}1$ MSE & $T{=}24$ MSE & $T{=}24$ $R^2$ & 参数量 \\
\midrule
DeReFusion（加法融合）      & 0.01647 & \textbf{0.06230} & \textbf{0.8691} & 28{,}436 \\
revin-DLinear（线性基座）   & \textbf{0.01577} & 0.07007 & 0.8528 & 4{,}664 \\
gatev1（波动率感知门控）    & ---     & 0.06767 & 0.8578 & 28{,}580 \\
gatev2（可学习门控）        & 0.02526 & 0.07191 & 0.8489 & $\sim$28k \\
\bottomrule
\end{tabular}
\end{table}

在 $T=24$ 上论文排序复现：加法融合的 MSE 比线性基座好 $11.1\%$，而两个门控变体都\textbf{劣于}无参数加法。
在 $T=1$ 上线性基座略优，与"先分解、再用简单模型"的长线性预测器行为一致。把同一协议扩展到另外两个资产
（$T=24$）：BTCUSD 上 DeReFusion $0.21797$ 对 DLinear $0.22322$（$+2.4\%$），
ETHUSD 上 $0.14602$ 对 $0.14876$（$+1.8\%$），即加法优势的量级排序为
GSPC $\gg$ BTCUSD $\approx$ ETHUSD。

\section{波动率工况分层}
\label{sec:regime}

每个测试窗口的工况由\textbf{仅使用输入窗口}的因果波动率度量确定：对数收益标准差，以及去趋势版本
$RV^{rel}_t = RV_t / \mathrm{median}(RV_{t-120:t-1})$。表~\ref{tab:regime} 报告非线性分支相对线性基座
在各分层的相对改善（主口径为相对波动率），以及交互效应 $\Delta_{high50} - \Delta_{low50}$。

\begin{table}[!t]
\centering
\small
\caption{分层表现（相对波动率）。$\Delta = $ MSE$_{非线性}-$MSE$_{线性}$，负值偏向非线性。
$^\dagger$ 表示 95\% 置信区间不含 0。}
\label{tab:regime}
\begin{tabular}{@{}lcccc@{}}
\toprule
资产 & 最平静 20\% & 高波动半区 & 最剧烈 20\% & 交互效应 [95\% CI] \\
\midrule
GSPC (2021) & $+6.0\%$      & $+14.2\%^\dagger$ & $+21.6\%^\dagger$ & $-0.01205\,[-0.01545,-0.00876]^\dagger$ \\
GSPC (2022) & $-2.7\%$      & $+23.2\%^\dagger$ & $+29.6\%^\dagger$ & $-0.02248\,[-0.02719,-0.01792]^\dagger$ \\
ETHUSD      & $+16.7\%^\dagger$ & $+6.3\%^\dagger$ & $+14.2\%^\dagger$ & $-0.00943\,[-0.01881,+0.00033]$ \\
BTCUSD      & $+3.5\%$      & $-1.0\%$          & $-3.5\%$          & $+0.01502\,[+0.00615,+0.02378]^\dagger$ \\
\bottomrule
\end{tabular}
\end{table}

由此得出两点。第一，标普 500 的效应\textbf{跨种子稳健}：两个种子都给出同量级且显著为负的交互效应，
而平静工况下两个分支无法区分。第二，BTCUSD 显著反号，因此"非线性优势随波动率上升"**不能**作为普适判据。
一项方法学检查解释了此前部分困惑：原始的波动率度量与时间强混淆（$RV$ 与样本序号的相关系数：
GSPC 为 $+0.44$，BTCUSD 为 $-0.67$）。去趋势后的 $RV^{rel}$ 消除了该趋势；在它之下，GSPC 中间五分位
出现的"线性反超"消失。因此报告中以 $RV^{rel}$ 为主口径。

\section{自适应融合（门控）没有收益}
\label{sec:gate}

对工况结构的自然利用方式是门控 $\alpha(X)$ 混合两个分支。框架本身恰好提供了该变体，我们在结构最强、
跨种子最稳健的资产上运行了它（GSPC，$T=24$）：波动率感知门控的 MSE 为 $0.06767$，即比无参数加法
（$0.06230$）\textbf{差 $8.6\%$}，仅比纯线性基座（$0.07007$）好 $3.4\%$。由于非线性分支在平静工况下
从不劣于线性分支，最优门控权重实际上恒等于 1，门控\textbf{没有可切换的空间}。结合此前可学习门控的负结果，
这一证据支持\textbf{停止自适应融合这条线}，而不是在其之上继续搭路由器。

\section{容量敏感性：此前的"刚性"部分是假象}
\label{sec:capacity}

下一节的代理实验用一个容量对齐的 MLP 与线性算子对比（$9{,}240$ 对 $9{,}431$ 个参数，$+2.1\%$）。
该对齐把 384 维输入压进 23 维隐层——这是严重的瓶颈。因此我们把隐层宽度作为\textbf{唯一变量}
（$23 \to 64 \to 128$；数据、特征、切分、优化器、学习率、批大小、训练轮数、早停、种子与评估协议全部不变）
重跑结构状态评估。表~\ref{tab:cap} 按宽度与资产给出 12 个结构状态 $\times$ 3 个种子的平均 $\Delta$MSE。

\begin{table}[!t]
\centering
\small
\caption{容量敏感性。平均 $\Delta$MSE（负值 $=$ 非线性更优）；括号内为 36 个（状态, 种子）观测中
偏好非线性的个数。}
\label{tab:cap}
\begin{tabular}{@{}lcccc@{}}
\toprule
隐层宽度 & MLP 参数量 & GSPC & BTCUSD & ETHUSD \\
\midrule
23  & 9{,}431  & $+0.0598$ (0/36)  & $+0.1783$ (0/36)  & $+0.0145$ (12/36) \\
64  & 26{,}200 & $+0.0250$ (7/36)  & $+0.1888$ (0/36)  & $\mathbf{-0.0281}$ (36/36) \\
128 & 52{,}376 & $+0.0061$ (8/36)  & $+0.2998$ (0/36)  & $\mathbf{-0.0321}$ (36/36) \\
\bottomrule
\end{tabular}
\end{table}

三项预先指定的检查给出一致读数。（i）\textbf{符号反转}：不同宽度之间共 25 例状态级反转，
其中 23 例属于 ETHUSD（三个种子皆有），8 例属于 GSPC（**全部且仅出现在 seed 2023**，即种子不稳定）；
BTCUSD 为零例。（ii）\textbf{状态间跨度}：$\Delta$MSE 的状态间极差随容量\textbf{缩小}
（GSPC $0.106 \to 0.060$，ETHUSD $0.089 \to 0.040$）；BTCUSD 的扩大来自单一高误差状态的量级放大。
（iii）\textbf{逐种子方向一致性}从 $31/36$ 降到 $29/36$、$28/36$。综合来看：容量充分后 ETHUSD 对非线性
算子的偏好是\textbf{资产级全体}的（十二个状态全部反转，而非某个子集），因此这是\textbf{资产级异质性}
而非\textbf{状态级专业化}。据此，早先"线性全面占优"的表述记为\textbf{部分由容量造成}，并且我们
\textbf{不}把容量扫描当作路由的证据。

\section{结构感知路由：最小受控代理实验}
\label{sec:routing}

\subsection{设计}
十二个结构状态在观察任何结果\textbf{之前}定义，使用因果窗口特征（滞后 1/5/10 的自相关、趋势 $t$ 统计量、
符号持续性、阈值固定为 $c=3\sigma$ 的跳变比例、去趋势相对波动率、偏度与峰度）。二值状态以训练集
中位数作为阈值；另外报告一个复合的\textbf{结构强度}状态与 $K{=}3$ 的 $k$-means 交叉检验。
两个算子分别为线性映射（$9{,}240$ 参数）与单隐层 23 单元的 MLP（$9{,}431$ 参数），
在完全一致的协议下训练（同种子、优化器、步数、批大小、学习率、早停与预处理）。

\subsection{结果}
在容量对齐的条件下，GSPC 与 BTCUSD 在\textbf{全部}十二个状态上都是线性算子更优（GSPC 的 $\Delta$MSE
介于 $+0.047$ 与 $+0.081$，三个种子全部显著；BTCUSD 介于 $+0.129$ 与 $+0.448$）。只有 ETHUSD 出现
两种符号的状态间差异（结构弱的状态偏向非线性，最多 $-0.013$；结构强的状态偏向线性，$+0.033$）。
没有任何状态的偏好能通过容量扫描的检验。\textbf{主要发现：在容量对齐的尺度上，未建立稳健的算子专业化。}
唯一跨资产一致的现象是：平静工况下两个算子无法区分。

\section{九个资产的资产级算子异质性}
\label{sec:assets}

最后，用同一冻结协议（$T=24$，seed 2021）在另外七个资产上运行，以刻画交互效应的\textbf{符号}。
GSPC 的两个种子合并为一个资产，故 $N=9$。表~\ref{tab:assets} 报告交互效应及其置信区间，
表~\ref{tab:assoc} 报告资产结构特征与交互效应之间的探索性关联，并含 leave-one-asset-out（LOO）检查。

\begin{table}[!t]
\centering
\small
\caption{资产级交互效应（相对波动率）。$^\dagger$ 表示置信区间不含 0。}
\label{tab:assets}
\begin{tabular}{@{}lccc@{}}
\toprule
资产 & $\Delta_{interaction}$ & 95\% CI & 是否显著 \\
\midrule
EURUSD & $-0.1371$ & $[-0.164,-0.111]$ & 是$^\dagger$ \\
GSPC   & $-0.0173$ & $[-0.027,-0.009]$ & 是$^\dagger$ \\
USDJPY & $-0.0122$ & $[-0.026,+0.002]$ & 否 \\
ETHUSD & $-0.0094$ & $[-0.019,+0.000]$ & 否 \\
DJI    & $-0.0084$ & $[-0.019,+0.003]$ & 否 \\
TM     & $-0.0081$ & $[-0.028,+0.011]$ & 否 \\
BTCUSD & $+0.0150$ & $[+0.006,+0.024]$ & 是$^\dagger$ \\
BABA   & $+0.0405$ & $[+0.022,+0.060]$ & 是$^\dagger$ \\
NVO    & $+0.0976$ & $[+0.046,+0.149]$ & 是$^\dagger$ \\
\bottomrule
\end{tabular}
\end{table}

\begin{table}[!t]
\centering
\small
\caption{资产结构特征与 $\Delta_{interaction}$ 的探索性 Spearman 关联（$N=9$）及 LOO 稳健性。}
\label{tab:assoc}
\begin{tabular}{@{}lcccc@{}}
\toprule
特征 & $\rho$ & $p$ & LOO 范围 & 是否稳健 \\
\midrule
$|\mathrm{ACF}_1|$      & $+0.683$ & $0.042$ & $[+0.548,+0.810]$ & 是 \\
已实现波动率            & $+0.633$ & $0.067$ & $[+0.476,+0.857]$ & 是 \\
峰度                    & $+0.533$ & $0.139$ & $[+0.333,+0.643]$ & 是 \\
趋势持续性              & $+0.100$ & $0.798$ & $[-0.286,+0.381]$ & 否（变号） \\
跳变比例                & ---      & ---     & --- & 恒定（0.0105） \\
\bottomrule
\end{tabular}
\end{table}

效应的符号\textbf{跨零}：六个资产为负、三个为正，其中五个统计显著。两个特征在 LOO 下稳定：
已实现波动率更高、一阶自相关更强的资产，在湍流工况下的非线性优势\textbf{更弱}（甚至反转）。
\textbf{这些关联是探索性的，仅作为候选规律报告}——九个资产下单个 $p$ 值不足以定论；跳变比例特征
在所有资产的中位数上恒定（每 95 个收益中 1 次超阈），即在预先固定的阈值下不带区分信息；
趋势持续性在 LOO 下变号，标记为\textbf{单资产敏感}。

\section{诊断与决策}
\label{sec:diagnosis}

\begin{table}[!t]
\centering
\small
\caption{五个诊断问题的汇总回答。}
\label{tab:q}
\begin{tabular}{@{}p{0.42\textwidth}p{0.48\textwidth}@{}}
\toprule
问题 & 回答 \\
\midrule
Q1 是否存在稳健的样本级路由证据？ &
否。代理实验与容量扫描都未给出稳定的"状态 $\rightarrow$ 算子"映射。 \\
Q2 是否存在真实的资产级算子异质性？ &
是。容量充分时 ETHUSD 在 36/36 个观测上偏好非线性；BTCUSD 在任何容量下偏好线性；
GSPC 接近打平。 \\
Q3 能否用现有结构特征解释？ &
仅能作为候选：已实现波动率与 $|\mathrm{ACF}_1|$ 与交互效应符号相关且通过 LOO；
但符号在九个资产上跨零。 \\
Q4 当前主要不确定性？ &
资产依赖与算子形态——不是结构表示，也不再是路由（路由已基本排除）。 \\
Q5 是否可重新考虑更强的非线性算子？ &
仅可作为候选之一，须与容量对齐的 MLP 对照，且不得由"ETHUSD 偏好非线性"外推到任一算子族。 \\
\bottomrule
\end{tabular}
\end{table}

本阶段记录的实践决策为：（i）\textbf{停止自适应融合}——波动率门控实测劣于无参数加法，且它想利用的
结构本身没有可切换的空间；（ii）\textbf{不构建神经路由器}——状态到算子的映射未被建立；
（iii）把算子选择当作\textbf{资产级配置问题}，这是证据支持的做法；（iv）任何结构化非线性算子都
必须以"可重复的环境"与"容量对齐的对比"为前提。

\section{局限}
\label{sec:limitations}

\textbf{算子保真度}：代理算子（展平窗口上的线性映射与 MLP）不同于框架内的分支——后者把 DLinear 分解
与 LSTM--Transformer 残差分支在 RevIN 之下组合。因此代理结果\textbf{不能}读作"非线性无用"；
在框架内，残差分支在标普 500 的 $T=24$ 上仍跨种子优于线性基座，两者并不矛盾。
\textbf{样本量}：资产级关联建立在九个资产、每资产一个种子（GSPC 两个）之上，关联系数的置信区间很宽，
分析为探索性。\textbf{工况定义}：工况为因果波动率的五分位；平静工况在算子间统计不可分，
因此"平静"是一个接受性结论而非独立效应。\textbf{环境}：全部为 CPU 运行，墙钟时间与论文 GPU 数值不可比，
仅解释相对比较。

\section{结论}
\label{sec:conclusion}

复现在 $T=24$ 上重现了已发表排序，且跨种子稳健。其后的诊断细化但不推翻该结果：非线性分支的边际价值是
\textbf{资产依赖}而非工况普适；波动率门控没有收益；容量对齐的对比表明，早先"线性一律占优"的偏好
\textbf{部分是容量瓶颈造成的假象}。我们\textbf{未}发现支持样本级路由的稳健实验依据，而确实发现了真实的
资产级算子异质性。因此证据支持的下一步问题不是"路由器应该如何选择"，而是
\textbf{"在哪些可重复的环境中，更强的非线性算子值得占用它的容量"}——且必须以容量对齐的基线为参照来回答。

\appendix
\section{可复现性说明}
\label{app:repro}

全部复现资产位于项目仓库的 \texttt{reproduction/} 下：分析脚本在
\texttt{reproduction/analysis/}，批次运行器在 \texttt{reproduction/batches/}，
各设置的结果在 \texttt{reproduction/results/}。代表性命令：

\begin{verbatim}
# 单次运行（CPU）
.venv\Scripts\python.exe run.py --task_name long_term_forecast --is_training 1 \
  --model_id GSPC_96_24 --model DeReFusion --data custom --root_path ./dataset/ \
  --data_path GSPC-2016-2025.csv --features MS --target Close --freq b \
  --seq_len 96 --label_len 48 --pred_len 24 --enc_in 4 --dec_in 4 --c_out 1 \
  --d_model 32 --moving_avg 25 --train_epochs 30 --batch_size 32 \
  --learning_rate 0.0001 --patience 5 --lradj cosine --rand_seed 2021 --no_use_gpu

# 工况分层（因果，主口径 = 相对波动率）
.venv\Scripts\python.exe reproduction/analysis/analyze_volatility_regimes.py \
  --csv dataset/GSPC-2016-2025.csv --tag GSPC --seed 2021 --rv-mode relative
\end{verbatim}

四条环境事实曾耗费实际时间，记录备用：数据管线无条件导入 \texttt{patool}、
\texttt{huggingface\_hub}、\texttt{sktime} 与 \texttt{datasets}；\texttt{joblib} 必须
pin 到 \texttt{1.5.3}；CPU 运行必须加 \texttt{--no\_use\_gpu}；\textbf{杀掉 DataLoader 工作进程
（\texttt{spawn\_main}）会让父训练进程死锁}，因此"卡死"必须靠\textbf{进程 CPU 时间}而不是日志输出来判断。

\begin{thebibliography}{9}
\bibitem{hsieh2026} Hsieh, C.-Y., Chen, Y.-C. (2026). DeReFusion: a decomposition-reconstruction
fusion framework for financial time series forecasting. \emph{Applied Soft Computing} 203, 116252.
\bibitem{zeng2023} Zeng, A., Chen, M., Zhang, L., Xu, Q. (2023). Are transformers effective for
time series forecasting? \emph{Proceedings of AAAI} 37(9), 11121--11128.
\end{thebibliography}

\end{document}
